\section{Compact mean families and the spectral gap for wave evolution}
\label{S:compact-mean-family}

Let $I = [0,1]$, let $K_x \subset \R^d$ be compact, and let $\sigma\mapsto
V_\sigma$ be continuous from $I$ into $C_c^\infty(K_x)$, real and nonnegative.
The index $\sigma$ is a family parameter, distinct from physical propagation
time $t$. Assume that one real $w \in H^1(\R^d)$ with $\|w\|_2 = 1$ satisfies
\begin{equation}\label{E:compact-common-trial}
  \|\nabla w\|_2^2-\int V_\sigma w^2\le-\varepsilon
  \quad\hbox{for all }\sigma\in I,\qquad \varepsilon>0.
\end{equation}
Put $M = \sup_I\|V_\sigma\|_\infty$ and $H_\sigma = -\Delta-V_\sigma$. The
phases $E,F$, inclusion $i$, data $e_0$ and observable $O_\phi$ are exactly
those of Section~\ref{S:constant-phase-young}.
Proposition~\ref{P:constant-compact-base}, with deterministic coefficient $h =
V_\sigma$ and noise amplitude zero, gives the ordinary physical flow
$P_\sigma(t)$ on $F$. Its $E$ restriction is $P_\sigma^E(t)$. Uniqueness of the
Duhamel solution makes finite-horizon constructions agree on overlaps; integer
horizons and autonomy give the same semigroup for every $t \ge 0$.

We import the following statements from the \emph{second edition}
of~\cite{teschl:14:mathematical}: Theorem~10.2, p.~258, gives the self-adjoint
$H^2$ operator, lower bound, essential spectrum $[0,\infty)$ and compact smooth
core for real compact smooth potentials in the dimensions considered here.
Theorem~10.13, p.~273, gives a simple positive ground state \emph{if the
spectral bottom is an eigenvalue}. Theorem~4.12 and Corollary~4.13,
pp.~140--141, give max--min and form monotonicity, with the essential threshold
used after discrete eigenvalues end and with the quadratic-form domain allowed.
Uniformity, fractional norms and all applications to $F$ are proved below.

\subsection{The ground family and the global phase}

\begin{proposition}[Ground states and a uniform rate gap]
  \label{P:compact-ground-family}
  The common form domain of $H_\sigma$ is $H^1$ with form $q_\sigma(u) =
  \|\nabla u\|_2^2-\int V_\sigma|u|^2$. Its bottom $E_1(\sigma)$ is a simple
  eigenvalue with a unique real positive-a.e. normalized ground state
  $p_\sigma$. Let $E_2(\sigma)$ be the second max--min level, set to zero when
  there is no second negative eigenvalue. Then
  \begin{align}
    -M&\le E_1(\sigma)\le-\varepsilon,\qquad
    E_1(\sigma)<E_2(\sigma)\le0,\label{E:compact-energy-levels}\\
    |E_j(\sigma)-E_j(\tau)|
    &\le\|V_\sigma-V_\tau\|_\infty,\qquad j=1,2.
    \label{E:compact-level-continuity}
  \end{align}
  The map $\sigma\mapsto p_\sigma$ is continuous into $H^2$ and bounded there.
  With $\lambda_\sigma = \sqrt{-E_1(\sigma)}$ and $\mu_\sigma =
  \sqrt{-E_2(\sigma)}$,
  \begin{equation}\label{E:compact-rate-gap}
    \sqrt\varepsilon\le\lambda_\sigma\le\sqrt M,\qquad
    d_0:=\inf_I(\lambda_\sigma-\mu_\sigma)>0.
  \end{equation}
\end{proposition}

\begin{proof}
  We first identify the common form domain and place the spectral bottom below
  the essential threshold. We then use simplicity and continuity to obtain a
  uniform gap, and compare ground states to prove their continuity in $H^2$.

  \medskip
  \noindent\textit{Step 1.~}We identify the common form domain.
  The imported Theorem~10.2 applies to $-V_\sigma$, which is real, compactly
  supported and in $L^2$ for $d \le 3$. It gives the stated self-adjoint
  operator on $H^2$ and essential threshold zero. Write $c = 1+M$ and $A_\sigma
  = H_\sigma+c \ge 1$. On $H^2$, integration by parts and the spectral theorem
  give
  \[
    \|A_\sigma^{1/2}u\|_2^2
    =q_\sigma(u)+c\|u\|_2^2,\qquad
    \|u\|_{H^1}^2\le\|A_\sigma^{1/2}u\|_2^2
    \le c\|u\|_{H^1}^2.
  \]
  Since $0 \le V_\sigma \le M$ and $c = 1+M$, the shifted potential term
  $c-V_\sigma$ lies between one and $c$. This gives the two form bounds on
  $H^2$. The issue is then to identify their completions, rather than assume
  the form domain from the formal expression.
  Spectral cutoffs $1_{[1,n]}(A_\sigma)u$ show that $H^2$ is a core for the
  half-power form norm. Fourier cutoffs show it is dense in $H^1$. The two
  equivalent completions are therefore $D(A_\sigma^{1/2})$ and $H^1$. Both norms
  dominate $L^2$ and give the same $L^2$ limit, so these are equal as subspaces,
  with the displayed form. This verifies the form domain used below.

  \medskip
  \noindent\textit{Step 2.~}We obtain the isolated simple ground state.
  The form lower bound and~\eqref{E:compact-common-trial} give $-M \le E_1 \le
  -\varepsilon$ by the imported form max--min principle. The bottom belongs to
  the spectrum and lies below the essential threshold, so it is an isolated
  eigenvalue of finite multiplicity. Only now apply the imported Theorem~10.13:
  it is simple with a positive ground state. Conjugation preserves the real
  operator; positivity and unit $L^2$ norm fix its real representative.

  \medskip
  \noindent\textit{Step 3.~}We make the spectral and rate gaps uniform.
  Simplicity, isolation and threshold zero give $E_1 < E_2 \le 0$. For every
  common unit form test, $|q_\sigma(u)-q_\tau(u)| \le
  \|V_\sigma-V_\tau\|_\infty$. Taking the same infima over orthogonal
  complements and suprema in the source max--min formula proves
  \eqref{E:compact-level-continuity}, including passage of $E_2$ to the
  threshold. Compactness gives $\Delta_*: = \min_I(E_2-E_1) > 0$. Thus,
  \[
    \lambda_\sigma-\mu_\sigma
    =\frac{E_2(\sigma)-E_1(\sigma)}
    {\lambda_\sigma+\mu_\sigma}
    \ge\frac{\Delta_*}{2\sqrt M}>0.
  \]
  The level difference is continuous and strictly positive at every point of
  the compact parameter interval. It therefore has positive minimum
  $\Delta_*$. Factoring $\lambda_\sigma^2-\mu_\sigma^2$ gives the displayed
  rate difference, whose denominator is at most $2\sqrt M$.
  Here $M \ge \varepsilon > 0$, and no division by $\mu_\sigma$ occurs.

  \medskip
  \noindent\textit{Step 4.~}We prove continuity of the positive normalized
  ground state.
  Put $d_V = \|V_\sigma-V_\tau\|_\infty$ and decompose $p_\tau = a p_\sigma+r$,
  with $r\perp p_\sigma$ and $0 < a \le 1$. The variational bound gives
  $q_\sigma(p_\tau) \le E_1(\sigma)+2d_V$, whereas the complementary spectral
  bound gives $q_\sigma(p_\tau) \ge E_1(\sigma)+\Delta_*\|r\|_2^2$. Hence,
  \begin{equation}\label{E:compact-ground-continuity}
    \|p_\tau-p_\sigma\|_2^2
    =2(1-a)\le2(1-a^2)\le4d_V/\Delta_*.
  \end{equation}
  Positivity of both ground states gives $a = \langle p_\tau,p_\sigma\rangle >
  0$. Orthogonality gives $a^2+\|r\|_2^2 = 1$. The upper form bound follows by
  changing the potential and then the lowest eigenvalue, each at cost at most
  $d_V$. The complementary lower bound uses the uniform gap on $r$. Combining
  them gives $\|r\|_2^2 \le 2d_V/\Delta_*$ and hence the displayed distance
  estimate.
  The eigenfunction equation is
  \[
    (1-\Delta)p_\sigma
    =(1-\lambda_\sigma^2+V_\sigma)p_\sigma.
  \]
  Its right side is continuous in $L^2$ and bounded there by $1+2M$. The inverse
  of $1-\Delta$ is an isometry from $L^2$ to the Bessel $H^2$ space, proving the
  asserted $H^2$ continuity and bound. This proves
  Proposition~\ref{P:compact-ground-family}.
\end{proof}

The rate gap~\eqref{E:compact-rate-gap} is a statement about the spectrum of
$H_\sigma$. The wave on $E$ is measured in the fixed norm of $E$, so we first
compare the two.

\begin{lemma}[Uniform fractional phase norms]\label{L:compact-fractional-phase}
  Let $B = 1-\Delta$, $c = 1+M$ and $A_\sigma = H_\sigma+c$. The quarter-power
  domain is $H^{1/2}$ and
  \begin{equation}\label{E:compact-quarter-power}
    \|B^{1/4}u\|_2\le\|A_\sigma^{1/4}u\|_2
    \le c^{1/4}\|B^{1/4}u\|_2.
  \end{equation}
  The norm
  \begin{equation}\label{E:compact-spectral-phase-norm}
    \|(u,z)\|_{E_\sigma}^2
    =\|A_\sigma^{1/4}u\|_2^2
    +\|A_\sigma^{-1/4}z\|_2^2
  \end{equation}
  is uniformly equivalent to the fixed norm on $E = H^{1/2}\times H^{-1/2}$,
  interpreting the negative entry by completion.
\end{lemma}

\begin{proof}
  We prove the needed fractional order comparison from resolvents. The
  resulting quarter-power norm comparison identifies the positive phase entry,
  and duality identifies the negative entry with uniform constants.

  \medskip
  \noindent\textit{Step 1.~}We compare the quarter-power norms and their
  domains.
  On the common form domain $H^1$, $B \le A_\sigma \le cB$ and all three
  operators are at least one. For positive self-adjoint $T,S$ with common form
  domain, $T \le S$ implies $(T+t)^{-1} \ge (S+t)^{-1}$ for $t > 0$. Indeed,
  \[
    \langle f,(T+t)^{-1}f\rangle
    =\sup_u\{2\operatorname{Re}\langle f,u\rangle
    -q_T(u)-t\|u\|_2^2\};
  \]
  completing the square in the closed form norm proves the identity, and
  comparison of the suprema proves the order. Consequently, $I-t(T+t)^{-1} \le
  I-t(S+t)^{-1}$. The scalar identity
  \[
    \sqrt x=\frac1\pi\int_0^\infty
    t^{-1/2}\frac{x}{x+t}\ud t
  \]
  follows by $t = xr^2$ and $\int_0^\infty(1+r^2)^{-1}\ud r = \pi/2$. Spectral
  integration and positive Tonelli, first over finite intervals and then by
  monotone convergence, give $T^{1/2} \le S^{1/2}$ as forms, including domain
  inclusions. Apply this to $B \le A_\sigma \le cB$ to obtain
  \eqref{E:compact-quarter-power}: the form of $A_\sigma^{1/2}$ is the squared
  norm of $A_\sigma^{1/4}$. The lower and upper comparisons give both domain
  inclusions, so they also give
  equality of the quarter-power domains with
  $H^{1/2}$.

  The supremum is over the common form domain. Increasing the quadratic form
  decreases every expression inside that supremum. This explains why the
  resolvent order is opposite to the original form order.

  \medskip
  \noindent\textit{Step 2.~}We identify the negative phase norm by duality.
  Under the $L^2$ pairing the dual norm of $\|A_\sigma^{1/4}u\|_2$ at $z \in
  L^2$ is $\|A_\sigma^{-1/4}z\|_2$. Cauchy--Schwarz gives the upper bound, and a
  scalar multiple of $A_\sigma^{-1/2}z$ attains it. Comparison
  of these dual norms preserves the comparison with reversed inequalities. The
  constants remain uniform because the positive norm constants depend only on
  $c$. Completion
  and~\eqref{E:compact-quarter-power} identify this negative space with
  $H^{-1/2}$, uniformly in $\sigma$. Adding the squared norms proves the claim.
  This proves Lemma~\ref{L:compact-fractional-phase}.
\end{proof}

\begin{proposition}[The complementary wave on $E$]\label{P:compact-E-wave}
  Put
  \[
    e_\sigma^E=(p_\sigma,\lambda_\sigma p_\sigma),\qquad
    \ell_\sigma^E(u,z)
    =\frac{\langle p_\sigma,u\rangle+
    \langle p_\sigma,z\rangle/\lambda_\sigma}{2},
    \qquad \Pi_\sigma^E=e_\sigma^E\ell_\sigma^E.
  \]
  These are bounded real objects, continuous in their respective norms, with
  $\ell_\sigma^E(e_\sigma^E) = 1$, and
  \begin{align}
    P_\sigma^E(t)e_\sigma^E
    &=e^{\lambda_\sigma t}e_\sigma^E,\qquad
    \ell_\sigma^EP_\sigma^E(t)
    =e^{\lambda_\sigma t}\ell_\sigma^E,\label{E:compact-E-eigen}\\
    \|P_\sigma^E(t)(I-\Pi_\sigma^E)\|_{\mathcal{B}(E)}
    &\le C(1+t)e^{\mu_\sigma t},\qquad t\ge0.
    \label{E:compact-E-complement}
  \end{align}
  The constant is uniform over $\sigma \in I$.
\end{proposition}

\begin{proof}
  We estimate the spectral wave matrix in the uniformly equivalent phase norm,
  separating the complementary spatial spectrum from the two ground-state
  branches. We then identify this spectral matrix with the ordinary physical
  flow.

  \medskip
  \noindent\textit{Step 1.~}We write the spectral wave and its complementary
  spatial bound.
  Define real spectral functions $C_t(h),S_t(h)$ by $\cos(t\sqrt h),\sin(t\sqrt
  h)/\sqrt h$ for $h > 0$, by $1,t$ at zero, and by
  $\cosh(t\sqrt{-h}),\sinh(t\sqrt{-h})/\sqrt{-h}$ for $h < 0$. The spectral wave
  matrix is
  \begin{equation}\label{E:compact-spectral-wave}
    \begin{pmatrix}
      C_t(H_\sigma)&S_t(H_\sigma)\\
      -H_\sigma S_t(H_\sigma)&C_t(H_\sigma)
    \end{pmatrix}.
  \end{equation}
  Real potentials, their resolvents and spectral resolutions commute with
  conjugation. All these real functions and positive $A_\sigma$ powers preserve
  the real phases.

  The spatial projection $Q_\sigma^{\mathrm{sp}} = I-|p_\sigma\rangle\langle
  p_\sigma|$ commutes with $H_\sigma,A_\sigma$ and is contractive on both
  entries of~\eqref{E:compact-spectral-phase-norm}. On its support $h \ge E_2 =
  -\mu_\sigma^2$, the phase isometry gives scalar matrix entries
  \[
    C_t(h),\quad \sqrt{h+c}\,S_t(h),\quad
    -hS_t(h)/\sqrt{h+c},\quad C_t(h).
  \]
  The isometry multiplies the displacement by $(h+c)^{1/4}$ and the velocity
  by $(h+c)^{-1/4}$. Conjugating the two off-diagonal entries by these weights
  produces the displayed factors $\sqrt{h+c}$ and its reciprocal. Thus, an
  ordinary scalar matrix bound gives an operator bound in the spectral phase
  norm.
  For $h \ge 0$, use $|S_t(h)| \le \min(t,h^{-1/2})$ and split $h \le 1$ from $h
  \ge 1$; the lower off-diagonal entry is at most one. For $-\mu_\sigma^2 \le h
  < 0$, use $|C_t(h)| \le e^{\mu_\sigma t}$ and $|S_t(h)| \le t e^{\mu_\sigma
  t}$, with $|h| \le M$ and $1 \le h+c \le c$. Thus, the complementary spatial
  matrix has norm at most $C(1+t)e^{\mu_\sigma t}$, including $\mu_\sigma = 0$.

  At nonnegative $h \le 1$, the upper off-diagonal entry is bounded by
  $\sqrt{1+c}\,t$. At $h \ge 1$, it is bounded by $\sqrt{1+c}$. The lower
  entry is bounded by $\sqrt h/\sqrt{h+c} \le 1$. These are uniform bounds on
  the nonnegative part of the spectrum.

  \medskip
  \noindent\textit{Step 2.~}We retain the decaying ground branch in the phase
  complement.
  The two spatial ground branches are $(p_\sigma,\pm\lambda_\sigma p_\sigma)$.
  Besides $\Pi_\sigma^E$, their projections include the decaying branch
  \[
    (p_\sigma,-\lambda_\sigma p_\sigma)
    \frac{\langle p_\sigma,u\rangle
    -\langle p_\sigma,z\rangle/\lambda_\sigma}{2}.
  \]
  Their sum is the spatial ground projection on both phase entries. Both
  projection norms are uniformly bounded in the spectral phase norm because
  $\sqrt\varepsilon \le \lambda_\sigma \le \sqrt M$ and $1 \le
  c-\lambda_\sigma^2 \le c$. The decaying branch has factor $e^{-\lambda_\sigma
  t}$ and remains in $I-\Pi_\sigma^E$. The fractional norm equivalence now
  proves the stated complement estimate for the spectral matrix.

  Removing the growing phase branch does not remove both copies of the spatial
  ground state. The displayed decaying projection accounts for the remaining
  copy. Its exponential factor is at most one for $t \ge 0$, and hence is
  bounded by the same complementary estimate. The uniform equivalence of phase
  norms then returns that estimate to the fixed space $E$.

  \medskip
  \noindent\textit{Step 3.~}We identify the spectral matrix with the ordinary
  flow.
  We identify this matrix with the physical flow. The same scalar bounds at
  graph levels $(A_\sigma u,A_\sigma^{1/2}z)$ give continuous solutions in
  $H^2\times H^1$. These graph norms are uniformly equivalent to the Sobolev
  norms: the $H^1$ case was proved above, and $A_\sigma-B = M-V_\sigma$ is
  bounded by $M$ while $A_\sigma,B \ge 1$. Spectral cutoffs give $u_t = z$ in
  $H^1$ and $z_t = -H_\sigma u$ in $L^2$, with continuous derivatives there.
  The
  free Fourier group is used at these graph levels so that the
  interaction-picture derivative is justified in $H^1\times L^2$. This
  derivative is not an operator-norm derivative on the whole phase. The
  free Fourier group on $H^1\times L^2$ has, after its Bessel isometry, the same
  bounded weighted entries and is strongly continuous by dominated Fourier
  convergence. Its generator is $(z,\Delta u)$ on $H^2\times H^1$, as the
  Fourier multipliers show. For compact smooth initial data, differentiating the
  free interaction picture at these graph levels gives Duhamel with force
  $(0,V_\sigma u)$. The identity also holds in $E$. Multiplication by $V_\sigma$
  maps $H^{1/2}$ into $L^2$ and then $H^{-1/2}$ with bound $M$, so Picard
  uniqueness identifies the spectral solution with the ordinary $E$ flow on
  these data. Density within $E$ and finite-horizon boundedness extend equality
  to all $E$. No density of $E$ in $F$ is used. Finally, the ground branches
  give~\eqref{E:compact-E-eigen}; their pairings are continuous because
  $p_\sigma \in H^2$. Continuity and uniform bounds of the displayed vectors and
  duals follow from Proposition~\ref{P:compact-ground-family}. This proves
  Proposition~\ref{P:compact-E-wave}.
\end{proof}

\subsection{The growing projection and its complementary wave evolution}

For the constant coordinates write $C(t)(c_1,c_2) = (c_1+tc_2,c_2)$ and
$D_\sigma(c_1,c_2) = (0,V_\sigma c_1) \in E$.
The data $e_0 = (0,1,0)$ lie entirely in these constant coordinates, outside
$iE$. So we follow the flow on all of $F$, where $D_\sigma$ couples the
constants to $E$.

\begin{proposition}[The growing projection on the wave phase space]
  \label{P:compact-F-projection}
  The physical flow satisfies
  \begin{equation}\label{E:compact-F-block}
    P_\sigma(t)(e,c)
    =\left(P_\sigma^E(t)e+
    \int_0^tP_\sigma^E(t-r)D_\sigma C(r)c\ud r,\ C(t)c\right).
  \end{equation}
  Set
  \begin{align}
    f_\sigma&=\frac{\int V_\sigma p_\sigma}{2\lambda_\sigma}>0,
    &\ell_\sigma^C(c)&=
    f_\sigma\left(\frac{c_1}{\lambda_\sigma}
    +\frac{c_2}{\lambda_\sigma^2}\right),
    \label{E:compact-constant-dual}\\
    e_\sigma&=ie_\sigma^E,
    &\ell_\sigma(e,c)&=\ell_\sigma^E(e)+\ell_\sigma^C(c),\qquad
    \Pi_\sigma=e_\sigma\ell_\sigma,\quad Q_\sigma=I-\Pi_\sigma.
    \label{E:compact-F-projection}
  \end{align}
  These vectors, duals and rank-one projections are norm-continuous and
  uniformly bounded, and
  \begin{equation}\label{E:compact-F-eigen}
    \ell_\sigma(e_\sigma)=1,\qquad
    P_\sigma(t)\Pi_\sigma=\Pi_\sigma P_\sigma(t)
    =e^{\lambda_\sigma t}\Pi_\sigma.
  \end{equation}
  For every fixed nonnegative smooth compact $\phi$ of integral one,
  \begin{equation}\label{E:compact-positive-margins}
    \inf_I\ell_\sigma(e_0)>0,\qquad
    \inf_I O_\phi(e_\sigma)>0.
  \end{equation}
\end{proposition}

\begin{proof}
  We derive the block formula by variation of constants in $E$. We then
  compute the constant-coordinate contribution to the growing dual, combine
  its finite-time contribution with an exponential tail, and obtain uniform
  positive margins by compactness.

  \medskip
  \noindent\textit{Step 1.~}We derive the block formula for the full phase.
  The constant coordinates obey $C(t)$ because the potential force has range in
  $iE$. To justify variation of constants on the $E$ component, let $g$ be a
  continuous $E$ force and substitute the ordinary free-Duhamel equation for
  $P_\sigma^E(t-r)g(r)$ in its time integral. With $B_\sigma^E(u,z) =
  (0,V_\sigma u)$, strong $E$-valued Fubini gives
  \begin{align*}
    &\int_0^t\int_r^t S(t-s)B_\sigma^E
    P_\sigma^E(s-r)g(r)\ud s\ud r\\
    &\qquad=\int_0^tS(t-s)B_\sigma^E
    \left[\int_0^sP_\sigma^E(s-r)g(r)\ud r\right]\ud s.
  \end{align*}
  The integration region is the triangle $0 \le r \le s \le t$. Reversing the
  order changes the outer variable from $r$ to $s$, with $0 \le r \le s$ in
  the inner integral. This is exactly the forcing term in the free-Duhamel
  equation for the proposed $E$ component.
  All norms are bounded on the finite triangle, so this exchange is absolute in
  $E$. Thus, $Y(t) = P_\sigma^E(t)e+\int_0^tP_\sigma^E(t-r)g(r)\ud r$ solves
  free Duhamel with force $B_\sigma^EY+g$. Take $g(r) = D_\sigma C(r)c$ and
  adjoin the constant block. This is the same $F$ Duhamel equation as the
  ordinary flow; uniqueness proves~\eqref{E:compact-F-block}. Its finite-horizon
  bounds are uniform since the family has bounded $L^\infty$ and $L^2$ norms on
  one compact support.

  \medskip
  \noindent\textit{Step 2.~}We compute the growing dual on the constant
  coordinates.
  The integral defining $f_\sigma$ is finite by $L^2$ Cauchy--Schwarz and
  positive because $V_\sigma$ is nonzero by the common trial and $p_\sigma > 0$
  a.e. Its coefficients are continuous and uniformly bounded by the $H^2$
  continuity and positive lower bound for the rate; no $p_\sigma \in L^1$
  assumption is used. Directly,
  \begin{align*}
    \ell_\sigma^ED_\sigma C(r)c&=f_\sigma(c_1+rc_2),\\
    \ell_\sigma^C(c)
    &=\int_0^\infty e^{-\lambda_\sigma r}
    f_\sigma(c_1+rc_2)\ud r,\\
    \ell_\sigma^C(C(t)c)
    &=\int_t^\infty e^{\lambda_\sigma(t-r)}
    f_\sigma(c_1+rc_2)\ud r.
  \end{align*}
  The first integral uses $\int_0^\infty e^{-\lambda_\sigma r}\ud r =
  1/\lambda_\sigma$ and $\int_0^\infty r e^{-\lambda_\sigma r}\ud r =
  1/\lambda_\sigma^2$. The last identity follows by translating the
  integration variable by $t$ and using $C(t)c = (c_1+tc_2,c_2)$.
  These integrals converge absolutely. Applying $\ell_\sigma$
  to~\eqref{E:compact-F-block}, the finite integral and this tail combine to
  give $\ell_\sigma P_\sigma(t) = e^{\lambda_\sigma t}\ell_\sigma$. The ground
  mode also gives $P_\sigma(t)e_\sigma = e^{\lambda_\sigma t}e_\sigma$ and
  $\ell_\sigma(e_\sigma) = 1$, proving~\eqref{E:compact-F-eigen}. Bounded
  bilinear multiplication and the preceding continuous coefficients prove norm
  continuity of all objects. Finally, $\ell_\sigma(e_0) =
  f_\sigma/\lambda_\sigma > 0$ and $O_\phi(e_\sigma) = \int\phi p_\sigma > 0$.
  Nonnegativity and positive integral of $\phi$ give a set of positive measure
  on which $\phi > 0$. Since $p_\sigma > 0$ almost everywhere, its pairing
  with $\phi$ is strictly positive. Continuity converts this pointwise
  positivity in the parameter into a positive uniform minimum.
  Both are continuous on compact $I$, proving
  \eqref{E:compact-positive-margins}. Positivity is in integrals, without a
  point trace or a one-dimensional ODE. This proves
  Proposition~\ref{P:compact-F-projection}.
\end{proof}

Proposition~\ref{P:compact-F-projection} isolates the growing part
$\Pi_\sigma$. It remains to bound the flow on the complement $Q_\sigma$. The
proof of Lemma~\ref{L:raw-block-indexed-cone} uses the normalized form of this
bound.

\begin{proposition}[Complementary wave evolution and parameter dependence]
  \label{P:compact-F-complement}
  Uniformly over $\sigma$ and $t \ge 0$,
  \begin{equation}\label{E:compact-F-complement}
    \|P_\sigma(t)Q_\sigma\|_{\mathcal{B}(F)}
    \le C(1+t)^3e^{\mu_\sigma t}.
  \end{equation}
  For every fixed $0 < \eta < d_0$,
  \begin{equation}\label{E:compact-normalized-complement}
    \|e^{-\lambda_\sigma t}P_\sigma(t)Q_\sigma\|
    \le C_\eta e^{-\eta t}.
  \end{equation}
  For every finite $T_0$,
  \begin{equation}\label{E:compact-parameter-flow}
    \sup_{0\le t\le T_0}\|P_\sigma(t)-P_\tau(t)\|_{\mathcal{B}(F)}
    \le C_{T_0}
    \big(\|V_\sigma-V_\tau\|_\infty+\|V_\sigma-V_\tau\|_2\big).
  \end{equation}
\end{proposition}

\begin{proof}
  We use the block formula to keep track of the growing response to the
  constant coordinates when subtracting the projection. We estimate the
  remaining integral and absorb its polynomial factor into the rate gap.
  Finally, we compare two potentials through their Duhamel equations.

  \medskip
  \noindent\textit{Step 1.~}We bound the complementary flow on the full phase.
  Write $Q_\sigma^E = I-\Pi_\sigma^E$. Split $D_\sigma$ into its growing and
  complementary parts in~\eqref{E:compact-F-block} and use the two tail
  integrals in the preceding proof. The exact $E$ component of
  $P_\sigma(t)Q_\sigma(e,c)$ is
  \begin{equation}\label{E:compact-complement-forcing}
    P_\sigma^E(t)Q_\sigma^Ee
    +\int_0^tP_\sigma^E(t-r)Q_\sigma^ED_\sigma C(r)c\ud r
    -e_\sigma^E\ell_\sigma^C(C(t)c),
  \end{equation}
  and the constant component is $C(t)c$. The final subtraction retains the
  growing response to constants in the projection. Since $\|Q_\sigma^ED_\sigma
  C(r)\| \le C(1+r)$, the integral in~\eqref{E:compact-complement-forcing} has
  norm at most
  \[
    Ce^{\mu_\sigma t}|c|
    \int_0^t(1+t-r)(1+r)\ud r
    =Ce^{\mu_\sigma t}(t+t^2+t^3/6)|c|.
  \]
  The integrand has one factor $1+t-r$ from the complementary $E$ flow and one
  factor $1+r$ from the free constant block. Its exponential is bounded by
  $e^{\mu_\sigma t}$ because $\mu_\sigma \ge 0$. Integrating the product of
  the two linear factors gives the displayed cubic polynomial.
  The first term is bounded by $C(1+t)e^{\mu_\sigma t}\|e\|_E$; the final term
  and constant component by $C(1+t)|c|$. Since $\mu_\sigma \ge 0$ this proves
  \eqref{E:compact-F-complement}. The boundedness of $(1+t)^3e^{-(d_0-\eta)t}$
  proves \eqref{E:compact-normalized-complement}. In particular, the
  constant-velocity polynomial is included in this estimate.

  After multiplying by $e^{-\lambda_\sigma t}$, the exponential factor is at
  most $e^{-d_0t}$. For $\eta < d_0$, the remaining factor
  $(1+t)^3e^{-(d_0-\eta)t}$ has finite supremum on $[0,\infty)$, giving the
  normalized estimate.

  \medskip
  \noindent\textit{Step 2.~}We prove finite-horizon dependence on the
  potential.
  Put $\mathcal{B}_\sigma = iB_F(V_\sigma)$ and $d_{\sigma,\tau} =
  \|V_\sigma-V_\tau\|_\infty+ \|V_\sigma-V_\tau\|_2$. The bound for $B_F$ gives
  $\|\mathcal{B}_\sigma-\mathcal{B}_\tau\| \le d_{\sigma,\tau}$. Subtract the
  two free-Duhamel equations on states. Their difference is the sum of the
  integral with $S_F(t-r)\mathcal{B}_\sigma(P_\sigma(r)-P_\tau(r))$ and that
  with $S_F(t-r)(\mathcal{B}_\sigma-\mathcal{B}_\tau)P_\tau(r)$. The uniform
  finite-horizon bounds control the second integral by a constant times
  $d_{\sigma,\tau}$. The first integral is bounded by a constant times the
  time integral of the unknown flow difference. Thus, the
  finite-horizon bounds and scalar Gronwall prove
  \eqref{E:compact-parameter-flow}. The operator norm is measurable:
  separability of $F$ expresses it as the supremum over a countable dense set of
  unit states, each with continuous image. Hence, Gronwall applies without a
  $\mathcal{B}(F)$-valued derivative or operator-norm continuity of the free
  group. This proves Proposition~\ref{P:compact-F-complement}.
\end{proof}

\subsection{The original nonnegative control and a common trial}

So far $V_\sigma$ was any continuous family
satisfying~\eqref{E:compact-common-trial}. We now produce such a family from a
nonnegative test $\Psi$ and the original covariance $\Gamma$, in two steps.

\begin{lemma}[The original mean and variance]\label{L:compact-raw-profile}
  Let $\Gamma \ne 0$ satisfy Hypothesis~\ref{H:homog}, let $a = \alpha_0$, and
  let real nonnegative nonzero $\Psi \in C_c^\infty((0,1)\times\R^d)$ be
  arbitrary, without a separability assumption. Define
  \begin{equation}\label{E:compact-raw-mean}
    h_\sigma(x)=\int_0^1|\sigma-u|^{-a}
    (\Gamma*\Psi(u,\cdot))(x)\ud u,\qquad
    \overline{q}(x)=\int_0^1\Psi(u,x)\ud u,\qquad Z=W(\Psi).
  \end{equation}
  Then $h$ is smooth in time and space on compact sets, and
  \begin{equation}\label{E:compact-raw-cost}
    0<v_\Psi:=\E Z^2
    =\int_0^1\int\Psi(\sigma,x)h_\sigma(x)\ud x\ud\sigma<\infty.
  \end{equation}
  Against every smooth compact positive-time test $F_0$,
  \[
    \E[W(\Psi)W(F_0)]=\int F_0(s,x)h_s(x)\ud x\ud s.
  \]
  Moreover, on $I$,
  \begin{equation}\label{E:compact-positive-mean-floor}
    h_\sigma\ge\Gamma*\overline{q},\qquad
    v_\Psi\ge Q(\overline{q})
    :=\int\Gamma(\ud z)\int\overline{q}(x)\overline{q}(x-z)\ud x>0.
  \end{equation}
  For every fixed mask, $V_\sigma = m_Rh_\sigma$ is a smooth compact family,
  Lipschitz in $\sigma$ in $L^\infty$ and $L^2$.
\end{lemma}

\begin{proof}
  We identify the original covariance mean and show its cost is finite.
  Positivity and homogeneity give a nonnegative, nonzero time-independent lower
  profile and a strictly positive covariance cost. We then differentiate the
  test, rather than the singular time kernel, to obtain regularity.

  \medskip
  \noindent\textit{Step 1.~}We identify the mean and bound its covariance
  cost.
  Choose $A \ge 1$ containing the spatial support of $\Psi$ in $B_A$. The
  original covariance identity and absolute Fubini give the stated mean pairing:
  the two compact tests restrict spatial differences to a compact set of finite
  $\Gamma$ mass, and the time singularity is locally integrable since $a < 1$.
  This identifies the linear covariance mean without evaluating noise at a
  point. The same calculation gives~\eqref{E:compact-raw-cost}. Its spatial
  integrand is bounded in integral by
  $\|\Psi\|_\infty^2|B_A|\Gamma(\overline{B}_{2A})$, and
  $\int_0^1\int_0^1|\sigma-u|^{-a}\ud u\ud\sigma = 2/((1-a)(2-a))$, proving
  finiteness.

  \medskip
  \noindent\textit{Step 2.~}We obtain strict positivity and the uniform mean
  floor.
  Nonnegative nonzero smooth $\Psi$ has positive integral, so $\overline{q} \ne
  0$. Nonzero $\Gamma$ charges some ball, and homogeneity gives $\Gamma(B_1) =
  r^{\alpha-d}\Gamma(B_r) > 0$. Thus, every origin-centered ball $B_r$, $r > 0$,
  has positive finite mass. The continuous compact autocorrelation
  $\int\overline{q}(x)\overline{q}(x-z)\ud x$ is positive at zero and bounded
  below on some ball. This proves $0 < Q(\overline{q}) < \infty$. Since
  $|\sigma-u|^{-a} \ge 1$ on the unit time square, positive Tonelli gives both
  inequalities in~\eqref{E:compact-positive-mean-floor}.

  The value of the spatial autocorrelation at zero is $\|\overline q\|_2^2 >
  0$.   The lower bound on the time kernel can then be integrated
  without changing signs, because both tests and the covariance measure are
  nonnegative.

  \medskip
  \noindent\textit{Step 3.~}We prove smoothness and the fixed-mask Lipschitz
  bounds.
  Extend $\Psi$ by zero in time and write
  \[
    h_s(x)=\int_\R |r|^{-a}
    \int\Gamma(\ud z)\Psi(s-r,x-z)\ud r.
  \]
  On compact $s,x$ sets the relevant $r,z$ range over fixed compact sets;
  $|r|^{-a}$ is integrable there. Every derivative falls on the smooth test, not
  on the singular time kernel. Dominated differentiation gives smoothness and
  continuity of all derivatives on compacts. Multiplication by $m_R$ now gives
  the claimed compact family and its fixed-$R$ Lipschitz bounds. This proves
  Lemma~\ref{L:compact-raw-profile}.
\end{proof}

Lemma~\ref{L:compact-raw-profile} gives a smooth compact family for every
fixed mask, but not yet the common negative
trial~\eqref{E:compact-common-trial}.

\begin{proposition}[Binding with the original control and cost fixed]
  \label{P:compact-original-common-trial}
  For every $\Psi$ in Lemma~\ref{L:compact-raw-profile}, a sufficiently large
  finite mask $m_R$ gives $V_\sigma = m_Rh_\sigma$ a common negative
  trial~\eqref{E:compact-common-trial}. Thus, all preceding ground, projection,
  positive-margin and complementary-wave conclusions apply. The rate is
  Lipschitz on $I$ and has $\lambda_0 > 0$. The original test $\Psi$, original
  variable $Z$ and variance $v_\Psi$ are unchanged by this choice. The mask can
  additionally cover any fixed compact variational test supports.
\end{proposition}

\begin{proof}
  We dilate a fixed trial function and choose the mask to contain its support.
  Homogeneity gives a potential contribution of order $L^{-\alpha}$, which
  dominates the kinetic contribution of order $L^{-2}$. We then apply the
  ground-family estimates at this fixed mask.

  \medskip
  \noindent\textit{Step 1.~}We bound the potential contribution of the dilated
  trial.
  Choose a real unit $L^2$ function $\psi \in C_c^\infty(B_2)$ with $\psi \ge
  c_0 > 0$ on $B_1$, and put $\psi_L(x) = L^{-d/2}\psi(x/L)$, $D_\psi =
  \|\nabla\psi\|_2^2$, where $L \ge \max(1,2A)$. Choose the even mask with $R
  \ge 2L$, so $V_\sigma = h_\sigma$ on $\operatorname{supp}\psi_L$. This is the
  allowed real mask $0 \le m_R \le 1$, equal to one on $B_R$ and with support in
  $\overline{B}_{2R}$. Positive Tonelli
  and~\eqref{E:compact-positive-mean-floor} give
  \begin{align}
    \int V_\sigma\psi_L^2
    &\ge c_0^2L^{-d}\int_{B_L}(\Gamma*\overline{q})(x)\ud x\nonumber\\
    &=c_0^2L^{-d}\int\overline{q}(y)\Gamma(B_L-y)\ud y\nonumber\\
    &\ge c_0^2L^{-d}\Gamma(B_{L/2})\int\overline{q}
    =C_\Gamma L^{-\alpha},\label{E:compact-trial-potential}\\
    C_\Gamma&=c_0^22^{\alpha-d}\Gamma(B_1)\int\overline{q}>0.\nonumber
  \end{align}
  On $B_L$, the dilated trial has squared value at least $c_0^2L^{-d}$. Every
  $y$ in the support of $\overline q$ has $|y| \le A$. Hence, translating
  $B_L$ by such a $y$ still contains $B_{L/2}$. Homogeneity changes its mass
  to $(L/2)^{d-\alpha}\Gamma(B_1)$, giving exactly the displayed constant.
  Here $B_{L/2} \subset B_L-y$ for $|y| \le A$ and $L \ge 2A$. All covariance
  masses in the argument are compact and finite. The kinetic energy is $D_\psi
  L^{-2}$. Since $\alpha < 2$, choose $L$ large enough that $\varepsilon_R: =
  C_\Gamma L^{-\alpha}-D_\psi L^{-2} > 0$. Then
  \[
    \|\nabla\psi_L\|_2^2-\int V_\sigma\psi_L^2
    \le-\varepsilon_R\quad\hbox{for every }\sigma\in I.
  \]

  \medskip
  \noindent\textit{Step 2.~}We fix a negative trial and obtain the rate
  regularity.
  Fix this finite $R$ before invoking Proposition~\ref{P:compact-ground-family}.
  Its level bound and the mean Lipschitz estimate give
  \[
    |\lambda_\sigma-\lambda_\tau|
    \le\frac{\|V_\sigma-V_\tau\|_\infty}{2\sqrt{\varepsilon_R}}
    \le C_R|\sigma-\tau|.
  \]
  The lowest eigenvalue remains at most $-\varepsilon_R$ for every parameter.
  Factoring the difference of its negative square roots gives a denominator
  $\lambda_\sigma+\lambda_\tau \ge 2\sqrt{\varepsilon_R}$. The level
  comparison then gives the displayed Lipschitz estimate.
  The positive mean integrates the entire time test, so $\lambda_0 > 0$ even
  though $\Psi(0,\cdot) = 0$. Only the trial and mask have grown; the original
  test, its variable and~\eqref{E:compact-raw-cost} were fixed first. Enlarging
  $R$ to contain additional fixed compact test supports preserves the trial and
  gives equality of masked and unmasked potential energies on those supports.
  This includes singular invariant-direction covariances and the one-dimensional
  white case. It is not an all-strength binding assertion for one fixed compact
  potential. Masks are fixed before source or time asymptotics; a later supremum
  over controls must follow their individual bounds. This proves
  Proposition~\ref{P:compact-original-common-trial}.
\end{proof}
